\section{Sixty-Six Settings, and the Five Mosaic Uses}
\label{sec:ch10-sixty-six-settings}
\index{router!configuration|(}%

Chapter~\ref{ch:how-a-federated-query-actually-runs} configured its router with
three environment variables and did not explain them. That is the whole of one
way to configure this thing. The other way is a file, and Mosaic uses it:

\begin{minted}{yaml}
version: "1"

execution_config:
  file:
    path: /etc/mosaic/supergraph.json
    watch: true
    watch_interval: 1s

listen_addr: "0.0.0.0:3002"
dev_mode: true
log_level: info
\end{minted}

\index{router!config.yaml@\code{config.yaml}}%
That is \code{router/config.yaml} at tag \code{ch10}, minus its comments. Five
settings. I counted the top-level keys the router accepts at 0.337.1 by reading
the configuration schema it publishes at the matching release tag, and there
are sixty-six of them.

\index{router!a misspelt key stops it}%
The root of that schema carries \code{additionalProperties: false}, and the
router enforces it rather than leaving it to your editor. I misspelt
\code{dev\_mode} as \code{dev\_moed} and the container exited 1 before serving
anything:

\begin{minted}[fontsize=\footnotesize,breaklines]{text}
Could not load config: errors while loading config files: router config validation error for /etc/mosaic-router/config.yaml: jsonschema validation failed
- at '': additional properties 'dev_moed' not allowed
\end{minted}

That is the right way round. A gateway which starts up carrying a setting you
thought you had set is worse than one that refuses to start at all.

The rest of this section is the two settings in that file whose behaviour is
not what the name suggests, and two more that Mosaic never sets and relies on
anyway.

\subsection*{The default listen address does not work in a container}

\index{router!listen\_addr@\code{listen\_addr}}%
\code{listen\_addr} defaults to \code{localhost:3002}. Inside a container that
is the container's own loopback interface, so the router starts, logs that it
is serving requests, reports the address it is listening on, and refuses every
connection from outside. Nothing about it looks broken from the log.

\code{0.0.0.0:3002} is the fix and it is one line. I am pointing at it because
the failure mode is a healthy-looking process that answers nothing, which is
the kind of thing you can lose twenty minutes to while checking your firewall.

\subsection*{The file beats the environment}

\index{router!configuration precedence}%
Here is the one that surprised me. Cosmo's documentation says
\enquote{Values specified in the config file have precedence over Environment
variables. This also includes empty values so only specify values that should be
overwritten}~\autocite{cosmo2026routerconfiguration}.

That is backwards from every container runtime I have used. The convention
everywhere else is that the image carries defaults and the deployment overrides
them with environment variables, because the environment is the thing you can
change without rebuilding.

I did not want to take a documentation sentence for it, and a precedence claim
needs a control or it proves nothing, so I ran two routers. Both in development
mode, both with \code{LOG\_LEVEL=debug} in the environment, differing only in
whether a configuration file said otherwise:

\begin{minted}[fontsize=\footnotesize]{text}
environment alone, LOG_LEVEL=debug:         9 debug lines
config.yaml says info, environment debug:   0 debug lines
\end{minted}

The file wins. The control is the first row: without it, zero debug lines in
the second run would be equally consistent with debug logging producing nothing
at startup, and I would have reported the finding either way.

\index{router!announces its own precedence}%
And the router does say so. It is the second line of
section~\ref{sec:ch10-one-process-two-files}'s startup log:

\begin{minted}[fontsize=\footnotesize,breaklines]{text}
INFO cmd/main.go:125 Config file provided. Values in the config file have higher priority than environment variables
\end{minted}

Which is fair of it, and does nobody any good. That line goes past at container
start, once, above two warnings that look more urgent. Where this bites is a
deployment that sets \code{LISTEN\_ADDR} or \code{LOG\_LEVEL} in its
orchestrator, against an image whose baked-in \code{config.yaml} happens to name
the same keys. The orchestrator is ignored and the manifest still says what you
meant.

My rule after this: put a key in one place. If the file names it, the
environment must not, and the way to make a value deployment-specific is
\code{"\$\{VAR\}"} interpolation inside the file rather than a second copy of
the setting outside it.

\subsection*{Localhost is not the container}

\index{router!localhost fallback inside Docker}%
The composed configuration from chapter~\ref{ch:composition} carries the
routing URLs it was given by the input file:

\begin{minted}{text}
catalog   http://localhost:5101/graphql
mosaic    http://localhost:5100/graphql
\end{minted}

Those are correct on the host and wrong inside a container, where
\code{localhost} is the router. By the argument above the storefront query
should have failed, and it did not, which meant I had a working system I could
not explain.

The answer is a setting nobody sets:

\begin{minted}{text}
localhost_fallback_inside_docker
\end{minted}

It defaults to \code{true}. When the router notices it is inside a container it
rewrites a \code{localhost} upstream to \code{host.docker.internal}, and its own
schema describes the setting as being for development and testing.

The control was that same key, set to \code{false} in the configuration file.
Same container, same graph, same query:

\begin{minted}[fontsize=\footnotesize]{json}
{"errors":[{"message":"Failed to fetch from Subgraph 'catalog'."}],"data":null}
\end{minted}

\index{errors!a failed fetch is an HTTP 200}%
Note the status code, which is 200. A GraphQL error is not a transport error,
and a router whose every upstream is unreachable answers 200 with nulls all
day. Anything watching status codes, a load balancer or a naive uptime check,
sees a healthy service. So does \code{/health/ready}, for the reason
section~\ref{sec:ch10-one-process-two-files} ended on: readiness is the router
reporting on itself. Chapter~\ref{ch:resilience-and-performance} is where
partial failure is the subject rather than an aside.

The message is also all you get. It names the subgraph and stops, because the
router's error propagation defaults to the mode called \code{wrapped}, which
exists to keep internal detail away from clients. At the default log level, with
development mode on, the router's own log adds that the request errored and not
why the dial failed.

\index{router!routing URL overrides}%
I have left the fallback on, and I want to be clear that this is a development
convenience rather than an architecture. It keeps one composed config working
whether the subgraphs were started with \code{dotnet run} or by compose, at the
cost of the router leaving the machine and coming back in through a published
port. The deployment answer is a key Mosaic does not set:

\begin{minted}{text}
overrides.subgraphs.<name>.routing_url
\end{minted}

That replaces a subgraph's URL at the router rather than at the composer, and
it is exactly the seam you want. The composed graph stays the same artefact in
every environment, and where each service actually answers becomes a property
of the environment. Mosaic does not need it yet.
Chapter~\ref{ch:ci-cd-for-a-federated-graph} will.

\subsection*{Reloading without restarting}

\index{router!hot reload}%
\code{execution\_config.file.watch} is the setting that makes
chapter~\ref{ch:composition} and this one join up. With it on, the router
re-reads the file when it changes and swaps the graph over.

I measured it rather than trusting the phrase \enquote{without downtime}.
\code{scripts/measure-router.mjs --reload} composes a second config from a
schema with one field removed, drops it over the watched file, and polls every
100 ms until the field disappears from introspection. At a one-second watch
interval the change appeared after 1{,}090 ms on one run and 1{,}326 ms on
another, and no request failed on either. The router logs the swap:

\begin{minted}[fontsize=\footnotesize,breaklines]{text}
INFO core/router.go:1744 Config file changed. Updating server with new config {"watcher_label": "execution_config", "path": "/etc/mosaic/supergraph.json"}
\end{minted}

There is an obvious way for this to fail that it turns out not to. The compose
file bind-mounts the single file rather than its directory, and a bind mount of
one file survives an overwrite but not a replacement: a tool that writes to a
temporary file and renames it over the target leaves the container holding the
old inode. So I checked what \code{wgc} does. Composing straight into the
mounted path leaves the inode unchanged, which the script asserts on every run.

\index{router!the watcher reads the timestamp}%
And a way for it to fail that I found by falling into it. The first version of
that script restored the healthy config with a file copy, and the router went
on serving the thin graph afterwards, for minutes. The watcher reads the file's
modification time, and on Windows a copy carries the source's timestamp across,
so the restored file looked older than the graph already loaded and no reload
fired. Rewriting the same bytes fixed it, and confirms what the watcher is
looking at: I rewrote a byte-identical file and the router reloaded two seconds
later. It notices a changed timestamp, not changed content.

Which means the loop the rest of this book will use is two commands and no
restart. Change a subgraph, export its schema, recompose, and the router in
front of it is serving the new graph a second later.

\index{router!configuration|)}%
